\centerline{\bf \Huge Текстовые задачи}
\centerline{\bf \Huge }

\begin{flushright}
        \scriptsize{\textbf{20:32. Прибыл Годжо Сатору}}
\end{flushright}

\centerline{\bf \Huge Задачи о числах}

\problem{1} Дано двузначное число. К нему приписали такое же число и получили четырехзначное число. Во сколько раз увеличилось при этом данное число?

\problem{2} Сумма двух чисел равна 231. Одно из них оканчивается нулем. Если этот нуль зачеркнуть, то получится второе число. Найти эти числа. А если сумма равна 462?

\problem{3} На сколько сумма всех четных чисел первой сотни больше суммы всех нечетных чисел этой сотни?

\centerline{\bf \Huge Задачи на движение}

\problem{4} Когда первый велосепедист проехал 1 км, вслед за ним выехал второй велосепедист, и с ним побежала собака, которая догнала первого велосепедиста и вернулась ко второму; так собака бегала между велосепедистами, пока второй догнал первого. Какое расстояние пробежала собака за это время? Известно, скорость собаки 500 м/мин, скорость первого велосипедиста 200 м/мин, второго - 250 м/мин.

\problem{5} Первую часть пути велосепедист-гонщик проехал со скоростью 40 км/ч, а вторую часть пути со скоростью 60 км/ч. Определите среднюю скорость гонщика на всем пути.

\problem{6} Половину пути мотоциклист ехал со скоростью 45 км/ч, затем задержался на 10 минут и чтобы наверстать время, он увеличил скорость на 15 км/ч. Каков весть путь?

\problem{7} Расстояние между двумя поселками равно 9 км. Дорога имеет подъем, равнинный участок и спуск. Скорость пешехода на подъеме равна 4 км/ч, на равнинном участке 5 км/ч, а на спуске 6 км/ч. Сколько километров состовляет равнинный участок, если пешеход проходит расстояние от одного поселка до другого и обратно за 3 часа 41 минуту?

